% chapters/chapter-01.tex — Inequalities I: AM–GM and Cauchy–Schwarz
% Demonstrates every component of the design system with real content.

\chapter{Inequalities I: \texorpdfstring{AM--GM}{AM-GM} and
         \texorpdfstring{Cauchy--Schwarz}{Cauchy-Schwarz}}
\label{ch:ineq1}

\begin{objectives}
\begin{itemize}
  \objective State and prove the arithmetic--geometric mean inequality
    for two and for $n$ positive reals, including Cauchy's
    forward--backward induction.
  \objective State and prove the Cauchy--Schwarz inequality in two
    ways (discriminant method and Lagrange identity).
  \objective Use the Engel (Titu) form of Cauchy--Schwarz to
    homogenize and crush fraction-sum inequalities.
  \objective Track \emph{equality cases} in every application---the
    difference between an answer and a solution.
\end{itemize}
\end{objectives}

\begin{prerequisites}
\begin{itemize}
  \item Fluency with algebraic manipulation: expanding, factoring,
    completing the square.
  \item Proof by induction (we will extend it in an unusual direction).
  \item Basic sums $\sumi$ notation and its manipulation.
\end{itemize}
\end{prerequisites}

\section{Why Inequalities?}

Exact equalities are the exception in mathematics; controlled
imprecision is the rule. An olympiad inequality problem asks you to
bound one quantity by another, and---crucially---to identify exactly
when the bound is achieved. That second demand is what turns
inequalities from cookbook algebra into proof-writing.

Almost everything in this chapter grows from a single observation:
squares are never negative. From $(a-b)^2 \ge 0$ we will build the
arithmetic--geometric mean inequality (\cref{sec:amgm}), and from a
cleverly chosen quadratic we will build Cauchy--Schwarz
(\cref{sec:cs}). Between them, these two results settle a remarkable
share of competition problems; the classic references are
\cite{steele2004, hlp1952}, and \cite{engel1998} catalogs them by
tactic.

\begin{keyidea}[One source, many rivers]\label{key:source}
Whenever you meet a symmetric inequality in positive variables, ask:
\emph{can I trace it back to a sum of squares?} The three master tools
of this chapter---AM--GM, Cauchy--Schwarz, and the Engel form---are
all squares in disguise.
\end{keyidea}

\section{The Arithmetic--Geometric Mean Inequality}
\label{sec:amgm}

\begin{definition}[Arithmetic and geometric means]\label{def:amgm}
For positive real numbers $a_1, a_2, \dots, a_n$, the
\emph{arithmetic mean} $A$ and \emph{geometric mean} $G$ are
\[
  A \coloneqq \frac{a_1 + a_2 + \cdots + a_n}{n},
  \qquad
  G \coloneqq \sqrt[n]{a_1 a_2 \cdots a_n}.
\]
\end{definition}

We begin with the case $n = 2$, where the entire content is one
perfect square.

\begin{theorem}[AM--GM, two variables]\label{thm:amgm2}
For all positive reals $a, b$,
\[
  \frac{a+b}{2} \;\ge\; \sqrt{ab},
  \qquad\text{with equality if and only if } a = b.
\]
\end{theorem}

\begin{proof}
For any real $t$ we have $(t)^2 \ge 0$. Taking $t = a - b$ and
expanding,
\[
  (a-b)^2 = a^2 - 2ab + b^2 \ge 0
  \quad\Longrightarrow\quad
  a^2 + b^2 \ge 2ab .
\]
Add $2ab$ to both sides:
\[
  a^2 + 2ab + b^2 \ge 4ab
  \quad\Longleftrightarrow\quad
  (a+b)^2 \ge 4ab .
\]
Both sides are positive, so taking square roots preserves the
inequality: $a + b \ge 2\sqrt{ab}$, which is the claim. Equality
forces $(a-b)^2 = 0$, i.e.\ $a = b$.
\end{proof}

\begin{figure}[ht]
  \centering
  \begin{tikzpicture}[scale=1.35, line cap=round]
    % Square of side a+b (a = 2.2, b = 1.0), split into a 2x2 grid:
    % top-left a^2, top-right ab, bottom-left ab, bottom-right b^2.
    \draw[fill=navylight]  (0,0)   rectangle (2.2,3.2);   % left column
    \draw[fill=goldlight]  (2.2,0) rectangle (3.2,3.2);   % right column
    \draw[fill=goldlight]  (0,0)   rectangle (2.2,1.0);   % bottom-left ab
    \draw[very thick, color=navy] (0,0) rectangle (3.2,3.2);
    % Partition lines: vertical after a=2.2, horizontal after b=1.0.
    \draw[thick, color=navy] (2.2,0) -- (2.2,3.2);
    \draw[thick, color=navy] (0,1.0) -- (3.2,1.0);
    % Area labels
    \node at (1.1,2.1)  {\large $a^2$};
    \node at (2.7,2.1)  {\large $ab$};
    \node at (1.1,0.5)  {\large $ab$};
    \node at (2.7,0.5)  {\large $b^2$};
    % Side length labels
    \node[below=2pt] at (1.1,0) {\small $a$};
    \node[below=2pt] at (2.7,0) {\small $b$};
    \node[left=2pt]  at (0,0.5) {\small $b$};
    \node[left=2pt]  at (0,2.1) {\small $a$};
    % Side brace: total side a+b
    \draw[decorate, decoration={brace, amplitude=5pt}, color=slate]
      (3.2,0) -- (3.2,3.2) node[midway, right=6pt, color=slate] {\small $a+b$};
  \end{tikzpicture}
  \caption{The proof of \cref{thm:amgm2} as geometry. The square of
  side $a+b$ has area $(a+b)^2 = a^2 + 2ab + b^2$; since $a^2 + b^2
  \ge 2ab$ (the two corner squares beat the two $ab$ rectangles in
  total), the area is always at least $4ab$. Equality is the perfectly
  symmetric picture $a = b$, where all four pieces are congruent.}
  \label{fig:amgm-square}
\end{figure}

The general case is one of the most elegant induction arguments in
all of mathematics---\emph{Cauchy's forward--backward induction}. The
statement $P(n)$ is: AM--GM holds for every $n$-tuple of positive
reals. We prove $P(2)$, then $P(n) \Rightarrow P(2n)$ (forward), then
$P(n) \Rightarrow P(n-1)$ (backward). From $P(2)$ we thus get $P(4),
P(8, \dots$ and backward reach \emph{every} integer.

\begin{theorem}[AM--GM, general]\label{thm:amgm}
For all positive reals $a_1, \dots, a_n$,
\[
  \frac{a_1 + \cdots + a_n}{n} \;\ge\; \sqrt[n]{a_1 \cdots a_n},
  \qquad\text{with equality iff } a_1 = \cdots = a_n.
\]
\end{theorem}

\begin{proof}
We follow the plan above.

\emph{Base $P(2)$:} \cref{thm:amgm2}.

\emph{Forward step $P(n) \Rightarrow P(2n)$:} Given $2n$ positive
reals $a_1, \dots, a_{2n}$, split them into two blocks of $n$. Writing
$A_1, G_1$ for the means of the first block and $A_2, G_2$ for the
second, and applying $P(n)$ to each block and then $P(2)$ to the pair
$(A_1, A_2)$:
\[
  \frac{a_1 + \cdots + a_{2n}}{2n}
  = \frac{A_1 + A_2}{2}
  \ge \frac{G_1 + G_2}{2}
  \ge \sqrt{G_1 G_2}
  = \sqrt[2n]{a_1 \cdots a_{2n}} .
\]
Equality forces $A_1 = G_1$, $A_2 = G_2$ (so each block is constant)
and $G_1 = G_2$ (so the constants agree): all $2n$ entries are equal.

\emph{Backward step $P(n) \Rightarrow P(n-1)$:} Given $n - 1$
positive reals $a_1, \dots, a_{n-1}$, let $m$ be their arithmetic
mean and \emph{adjoin the mean itself as an $n$-th entry}. The
enlarged list $a_1, \dots, a_{n-1}, m$ has arithmetic mean
\[
  \frac{(a_1 + \cdots + a_{n-1}) + m}{n}
  = \frac{(n-1)m + m}{n} = m,
\]
so applying $P(n)$ to it yields
\[
  m \;\ge\; \sqrt[n]{a_1 \cdots a_{n-1} \cdot m}
  \quad\Longrightarrow\quad
  m^n \ge a_1 \cdots a_{n-1} \cdot m
  \quad\Longrightarrow\quad
  m^{n-1} \ge a_1 \cdots a_{n-1},
\]
where we divided by $m > 0$. Taking $(n-1)$-th roots gives $m \ge
\sqrt[n-1]{a_1 \cdots a_{n-1}}$, which is $P(n-1)$. Equality in
$P(n)$ forces the enlarged list to be constant; since its last entry
is $m$, the original $n-1$ entries all equal $m$.

Starting from $P(2)$, forward induction gives $P(2^k)$ for all $k$;
for any $n$, choose $k$ with $2^k \ge n$ and descend by the backward
step. \end{proof}

\begin{example}[The workhorse $x + 1/x$]\label{ex:reciprocal}
For every $x > 0$,
\[
  x + \frac{1}{x} \;\ge\; 2\sqrt{x \cdot \tfrac{1}{x}} = 2,
  \qquad\text{equality iff } x = 1 .
\]
\Cref{fig:x-reciprocal} shows why this single line solves many
problems: the function dips to its minimum exactly once, tangentially.
\end{example}

\begin{figure}[ht]
  \centering
  \begin{tikzpicture}
    \begin{axis}[
      width=0.72\textwidth, height=5.6cm,
      axis lines = middle,
      xlabel={$x$}, ylabel={$y$},
      xmin=0, xmax=4.4, ymin=0, ymax=6.4,
      xtick={1,2,3,4}, ytick={2,4,6},
      tick label style={font=\small},
      label style={font=\small},
      grid=both, grid style={rulegray!35},
      major grid style={rulegray!55},
    ]
      \addplot[very thick, color=navy, domain=0.28:4.2, samples=140]
        {x + 1/x};
      \addplot[dashed, color=gold, thick] (0,2) -- (4.4,2);
      \addplot[dashed, color=slate, thick] (1,0) -- (1,2);
      \addplot[only marks, mark=*, mark size=2.2pt, color=gold]
        coordinates {(1,2)};
      \node[font=\small, color=slate, anchor=west] at (axis cs:2.15,1.4)
        {minimum $=2$ at $x=1$};
    \end{axis}
  \end{tikzpicture}
  \caption{$f(x) = x + 1/x$ lies above the line $y = 2$ for all
  $x>0$, touching it only at $x=1$: AM--GM made visible.}
  \label{fig:x-reciprocal}
\end{figure}

\begin{example}[Fixed perimeter, maximal area]\label{ex:perimeter}
Among all rectangles with perimeter $P$, the square has the largest
area.

\begin{exsolution}
Let the side lengths be $a, b > 0$ with $a + b = P/2$. By
\cref{thm:amgm2},
\[
  ab \;\le\; \left(\frac{a+b}{2}\right)^{\!2} = \frac{P^2}{16},
  \qquad\text{equality iff } a = b = \frac{P}{4}.
\]
The maximal area $P^2/16$ is achieved precisely by the square.
\end{exsolution}
\end{example}

\begin{pitfall}[Forgetting the domain]\label{pit:domain}
AM--GM demands \emph{positive} reals. The step $x^2 \le y^2
\Rightarrow x \le y$ used in \cref{thm:amgm2} fails for negative
values, and $\sqrt[n]{\cdot}$ is not even defined (in $\R$) for
negative inputs and even $n$. When a variable can vanish, run AM--GM
on $\varepsilon$-shifted variables or split off the degenerate case.
\end{pitfall}

\section{Cauchy--Schwarz and Its Engel Form}
\label{sec:cs}

\begin{theorem}[Cauchy--Schwarz]\label{thm:cs}
For all real numbers $a_1, \dots, a_n$ and $b_1, \dots, b_n$,
\[
  \left( \sumk a_k b_k \right)^{\!2}
  \;\le\;
  \left( \sumk a_k^2 \right)
  \left( \sumk b_k^2 \right),
\]
with equality iff the vectors $(a_k)$ and $(b_k)$ are proportional.
\end{theorem}

\begin{proof}[First proof (the discriminant trick)]
If all $a_k$ vanish the claim is trivial, so assume
$\sumk a_k^2 > 0$. For every real $t$,
\[
  f(t) \;\coloneqq\; \sumk (a_k t + b_k)^2
  = \Big(\sumk a_k^2\Big) t^2
  + 2\Big(\sumk a_k b_k\Big) t
  + \sumk b_k^2
  \;\ge\; 0,
\]
being a sum of squares. Thus the quadratic $f$ never dips below zero,
so its discriminant is at most $0$:
\[
  \Delta = 4\left(\sumk a_k b_k\right)^{\!2}
  - 4\left(\sumk a_k^2\right)\left(\sumk b_k^2\right) \le 0,
\]
which rearranges to the claim. Equality means $f$ has a real double
root $t_0$: then $a_k t_0 + b_k = 0$ for every $k$, i.e.\ the vectors
are proportional.
\end{proof}

\begin{proof}[Second proof (Lagrange identity)]
Expanding the right-hand side minus the left:
\[
  \left(\sumk a_k^2\right)\!\left(\sumk b_k^2\right)
  - \left(\sumk a_k b_k\right)^{\!2}
  = \sum_{1 \le i < j \le n} \left( a_i b_j - a_j b_i \right)^2 \!\ge 0 .
\]
Every summand on the right is a square. Equality forces
$a_i b_j = a_j b_i$ for all pairs $i<j$, which for nonzero vectors
says exactly: proportional.
\end{proof}

\begin{corollary}[Engel form / Titu's lemma]\label{cor:engel}
For positive reals $b_1, \dots, b_n$ and arbitrary reals
$a_1, \dots, a_n$,
\[
  \sumk \frac{a_k^2}{b_k}
  \;\ge\;
  \frac{(a_1 + \cdots + a_n)^2}{b_1 + \cdots + b_n},
\]
with equality iff the ratios $a_k / b_k$ are all equal.
\end{corollary}

\begin{proof}
Apply \cref{thm:cs} to the vectors
$\big(a_1/\sqrt{b_1}, \dots, a_n/\sqrt{b_n}\big)$ and
$\big(\sqrt{b_1}, \dots, \sqrt{b_n}\big)$:
\[
  (a_1 + \cdots + a_n)^2
  = \left(\sumk \frac{a_k}{\sqrt{b_k}} \cdot \sqrt{b_k}\right)^{\!2}
  \le \left(\sumk \frac{a_k^2}{b_k}\right)
       \left(\sumk b_k\right). \qedhere
\]
\end{proof}

The Engel form turns sums of fractions into one fraction---a
\emph{homogenizing} move of enormous competitive value.

\begin{example}[Nesbitt's inequality]\label{ex:nesbitt}
For $a, b, c > 0$,
\[
  \frac{a}{b+c} + \frac{b}{c+a} + \frac{c}{a+b} \;\ge\; \frac{3}{2}.
\]

\begin{exsolution}
Write each term with a square on top,
$\dfrac{a}{b+c} = \dfrac{a^2}{ab+ac}$, and apply \cref{cor:engel} to
the three fractions:
\[
  \frac{a^2}{ab+ac} + \frac{b^2}{bc+ba} + \frac{c^2}{ca+cb}
  \;\ge\;
  \frac{(a+b+c)^2}{2(ab+bc+ca)} .
\]
Now $(a+b+c)^2 \ge 3(ab+bc+ca)$, because the difference is
$\tfrac{1}{2}\big[(a-b)^2 + (b-c)^2 + (c-a)^2\big] \ge 0$. Hence the
right side is at least
$\tfrac{3(ab+bc+ca)}{2(ab+bc+ca)} = \tfrac{3}{2}$. Equality holds
exactly at $a = b = c$.
\end{exsolution}
\end{example}

\begin{strategy}[The Engel playbook]\label{str:engel}
When a target inequality involves $\sum \frac{(\text{something})^2}
{\text{something}}$---or can be massaged into that shape:
\begin{enumerate}[label=\arabic*., itemsep=2pt]
  \item force a square into each numerator
        ($\frac{a}{b+c} \mapsto \frac{a^2}{ab+ac}$);
  \item apply \cref{cor:engel} to collapse the sum to one fraction;
  \item clean up with whatever symmetric bound remains
        (often $(a+b+c)^2 \ge 3(ab+bc+ca)$);
  \item \emph{verify the equality case} before declaring victory.
\end{enumerate}
Many further worked examples of this playbook appear in
\cite{andreescu2000, kedlaya1999}.
\end{strategy}

\ifstudent\else
\begin{instructornote}
When teaching \cref{ex:nesbitt}, resist writing the Engel application
immediately. First let students try common denominators and watch the
combinatorial explosion---the motivation for homogenization lands much
harder after twenty seconds of suffering.
\end{instructornote}
\fi

\section{Problems}
\label{sec:problems-1}

\noindent Difficulty is marked on the five-star scale explained in the
front matter. Attempt each problem before reading its hint.

% exercises/problems-01.tex — Problems for Chapter 1
% Usage: \begin{exercise}[label=ex:1-N]{difficulty}
%        difficulty: 1=* routine ... 5=***** IMO P6 caliber

\begin{exercise}[label=ex:1-1]{1}
Prove that for all $x > 0$,
\[
  x + \frac{4}{x} \;\ge\; 4,
\]
and determine exactly when equality holds.
\end{exercise}

\begin{exercise}[label=ex:1-2]{1}
Prove that for all $a, b > 0$,
\[
  (a + b)\left( \frac{1}{a} + \frac{1}{b} \right) \ge 4 .
\]
\end{exercise}

\begin{exercise}[label=ex:1-3]{2}
Let $a, b, c$ be real numbers with $a + b + c = 1$. Prove that
\[
  a^2 + b^2 + c^2 \;\ge\; \frac{1}{3}.
\]
\end{exercise}

\begin{exercise}[label=ex:1-4]{2}
Prove that for all real $a, b, c$,
\[
  a^2 + b^2 + c^2 \;\ge\; ab + bc + ca .
\]
\end{exercise}

\begin{exercise}[label=ex:1-5]{2}
Let $x, y, z$ be positive reals with $xyz = 1$. Prove that
\[
  x + y + z \;\ge\; 3 .
\]
\end{exercise}

\begin{exercise}[label=ex:1-6]{3}
Let $x, y, z$ be positive reals. Prove that
\[
  (x + y)(y + z)(z + x) \;\ge\; 8xyz .
\]
\end{exercise}

\begin{exercise}[label=ex:1-7]{3}
Let $a, b, c$ be nonnegative reals with $a + b + c = 6$. Find the
maximum value of $\sqrt{a} + \sqrt{b} + \sqrt{c}$ and all triples
attaining it.
\end{exercise}

\begin{exercise}[label=ex:1-8]{3}
Let $x, y, z$ be positive reals. Prove that
\[
  \frac{x^2}{y+z} + \frac{y^2}{z+x} + \frac{z^2}{x+y}
  \;\ge\; \frac{x+y+z}{2}.
\]
\end{exercise}

\begin{challenge}[label=ch:1-1]{4}
Let $a, b, c$ be positive reals. Prove that
\[
  \frac{a^3}{a^2 + ab + b^2}
  + \frac{b^3}{b^2 + bc + c^2}
  + \frac{c^3}{c^2 + ca + a^2}
  \;\ge\; \frac{a+b+c}{3}.
\]
\source{Training problem (Shapiro--type)}
\end{challenge}

\begin{challenge}[label=ch:1-2]{5}
Let $a, b, c$ be positive reals. Prove that
\[
  \frac{a}{\sqrt{a^2 + 8bc}}
  + \frac{b}{\sqrt{b^2 + 8ca}}
  + \frac{c}{\sqrt{c^2 + 8ab}}
  \;\ge\; 1 .
\]
\source{IMO 2001, Problem 2}
\end{challenge}


\section{Hints}
\label{sec:hints-1}

% exercises/hints-01.tex — Hints for Chapter 1 problems
% Convention: one hint per problem, escalating specificity. Numbering
% matches the exercise counter; challenges are prefixed C.

\begin{itemize}[leftmargin=2.2em, itemsep=5pt, label={}]
  \item[\textbf{1.}] Direct AM--GM on the pair $\big(x, \tfrac{4}{x}\big)$.
    What is the product of the two terms?
  \item[\textbf{2.}] Expand, then pair each fraction with its
    reciprocal twin: $\tfrac{a}{b} + \tfrac{b}{a} \ge 2$.
  \item[\textbf{3.}] Cauchy--Schwarz with
    $(1, 1, 1) \cdot (a, b, c)$: bound $(a+b+c)^2$ from above by
    $3(a^2+b^2+c^2)$.
  \item[\textbf{4.}] Multiply the desired inequality by $2$ and
    recognize three perfect squares.
  \item[\textbf{5.}] AM--GM on three terms; the constraint $xyz = 1$
    makes the geometric mean equal to $1$.
  \item[\textbf{6.}] Bound each factor separately by AM--GM, then
    multiply the three bounds.
  \item[\textbf{7.}] Cauchy--Schwarz on
    $(\sqrt a + \sqrt b + \sqrt c)^2$; equality analysis will
    identify the maximizer.
  \item[\textbf{8.}] Engel form (\cref{cor:engel}) collapses the
    three fractions; the denominator becomes $2(x+y+z)$.
  \item[\textbf{C1.}] Prove the pointwise bound
    $\dfrac{a^3}{a^2+ab+b^2} \ge \dfrac{2a-b}{3}$
    by clearing denominators, then sum.
  \item[\textbf{C2.}] Engel form rewrites the left side as
    $\sum \frac{a^2}{a\sqrt{a^2+8bc}}$. Then Cauchy--Schwarz once
    more, expand $(a+b+c)^3$, and finish with six-term AM--GM.
\end{itemize}


\ifstudent\else
\section{Solutions}
\label{sec:solutions-1}
% exercises/solutions-01.tex — Full solutions for Chapter 1
% Rendered in the instructor edition only (guarded in the chapter).

% ---------- Exercise 1.1 -------------------------------------------------
\begin{solution}{ex:1-1}
Apply AM--GM (\cref{thm:amgm2}) to the pair
$\big(x, \tfrac{4}{x}\big)$, whose product is $4$:
\[
  x + \frac{4}{x} \;\ge\; 2\sqrt{x \cdot \frac{4}{x}} = 2\sqrt4 = 4 .
\]
Equality holds iff $x = \tfrac{4}{x}$, i.e.\ $x = 2$.
\answer{$x + 4/x = 4$ exactly at $x = 2$.}
\end{solution}

% ---------- Exercise 1.2 -------------------------------------------------
\begin{solution}{ex:1-2}
Expanding,
\[
  (a+b)\left(\frac1a + \frac1b\right)
  = 1 + \frac{a}{b} + \frac{b}{a} + 1
  = 2 + \frac{a}{b} + \frac{b}{a}.
\]
By AM--GM on the pair $\big(\tfrac{a}{b}, \tfrac{b}{a}\big)$, whose
product is $1$,
\[
  \frac{a}{b} + \frac{b}{a} \ge 2 ,
\]
so the original expression is at least $4$. Equality holds iff
$\tfrac{a}{b} = \tfrac{b}{a}$, i.e.\ $a = b$.
\end{solution}

% ---------- Exercise 1.3 -------------------------------------------------
\begin{solution}{ex:1-3}
By Cauchy--Schwarz (\cref{thm:cs}) applied to $(1,1,1)$ and
$(a,b,c)$,
\[
  (a+b+c)^2 \;\le\; 3\,(a^2+b^2+c^2).
\]
Since $a+b+c = 1$, this gives $a^2+b^2+c^2 \ge \tfrac13$. Equality
holds iff $(a,b,c)$ is proportional to $(1,1,1)$, i.e.\
$a=b=c=\tfrac13$.
\end{solution}

% ---------- Exercise 1.4 -------------------------------------------------
\begin{solution}{ex:1-4}
Multiply the target by $2$ and rearrange:
\[
  2(a^2+b^2+c^2) - 2(ab+bc+ca)
  = (a-b)^2 + (b-c)^2 + (c-a)^2 \;\ge\; 0,
\]
a sum of squares. Equality holds iff $a = b = c$. Note this is
exactly the identity used in \cref{ex:nesbitt}.
\end{solution}

% ---------- Exercise 1.5 -------------------------------------------------
\begin{solution}{ex:1-5}
By AM--GM for three variables (\cref{thm:amgm}, $n=3$),
\[
  \frac{x+y+z}{3} \;\ge\; \sqrt[3]{xyz} = \sqrt[3]{1} = 1,
\]
so $x+y+z \ge 3$. Equality holds iff $x=y=z=1$.
\end{solution}

% ---------- Exercise 1.6 -------------------------------------------------
\begin{solution}{ex:1-6}
Apply \cref{thm:amgm2} to each factor:
\[
  x+y \ge 2\sqrt{xy}, \qquad
  y+z \ge 2\sqrt{yz}, \qquad
  z+x \ge 2\sqrt{zx}.
\]
All three left sides are positive, so multiplying the inequalities is
legitimate:
\[
  (x+y)(y+z)(z+x) \;\ge\; 8\sqrt{xy \cdot yz \cdot zx}
  = 8\sqrt{x^2y^2z^2} = 8xyz .
\]
Equality holds iff $x=y=z$.
\end{solution}

% ---------- Exercise 1.7 -------------------------------------------------
\begin{solution}{ex:1-7}
By Cauchy--Schwarz applied to
$(\sqrt a, \sqrt b, \sqrt c)$ and $(1,1,1)$,
\[
  (\sqrt a + \sqrt b + \sqrt c)^2
  \;\le\; 3\,(a+b+c) = 18,
\]
so $\sqrt a + \sqrt b + \sqrt c \le 3\sqrt2$.
\answer{Maximum $= 3\sqrt2$, attained only at $(a,b,c) = (2,2,2)$.}
Equality in Cauchy--Schwarz requires
$(\sqrt a, \sqrt b, \sqrt c)$ proportional to $(1,1,1)$, i.e.\
$a=b=c$; combined with $a+b+c=6$ this forces $a=b=c=2$, which
indeed gives $3\sqrt2$. Since the domain $\{a,b,c \ge 0,\;
a+b+c=6\}$ is compact and the objective is continuous, the maximum
exists, so the bound is sharp.
\end{solution}

% ---------- Exercise 1.8 -------------------------------------------------
\begin{solution}{ex:1-8}
Apply the Engel form (\cref{cor:engel}) to the three fractions:
\[
  \frac{x^2}{y+z} + \frac{y^2}{z+x} + \frac{z^2}{x+y}
  \;\ge\;
  \frac{(x+y+z)^2}{(y+z)+(z+x)+(x+y)}
  = \frac{(x+y+z)^2}{2(x+y+z)}
  = \frac{x+y+z}{2} .
\]
Equality holds iff $\tfrac{x}{y+z} = \tfrac{y}{z+x} = \tfrac{z}{x+y}$,
which (cross-multiplying and subtracting) forces $x=y=z$.
\end{solution}

% ---------- Challenge 1 ---------------------------------------------------
\begin{solution}{ch:1-1}
Following the hint, we prove the pointwise bound
\[
  \frac{a^3}{a^2+ab+b^2} \;\ge\; \frac{2a-b}{3} .
\]
Since $a^2+ab+b^2 > 0$, clearing denominators shows the difference is
\[
  3a^3 - (2a-b)(a^2+ab+b^2)
  \;=\; 3a^3 - \big(2a^3 + a^2b + ab^2 - b^3\big)
  \;=\; (a+b)(a-b)^2 \;\ge\; 0,
\]
where the factorization is verified by expanding
$(a+b)(a-b)^2 = a^3 - a^2b - ab^2 + b^3$. Equality holds iff $a=b$.
Applying the same bound to the other two cyclic terms and summing,
\[
  \sum_{\text{cyc}} \frac{a^3}{a^2+ab+b^2}
  \;\ge\;
  \frac{(2a-b) + (2b-c) + (2c-a)}{3}
  = \frac{a+b+c}{3} .
\]
Equality requires $a=b$, $b=c$, and $c=a$ simultaneously---that is,
$a=b=c$, where both sides indeed agree (each fraction equals
$\tfrac{a}{3}$).
\end{solution}

% ---------- Challenge 2 (IMO 2001/2) --------------------------------------
\begin{solution}{ch:1-2}
Write $S = a+b+c$. We chain three estimates.

\emph{Step 1 (Engel).} Writing each term as a square over a product,
\[
  \sum_{\text{cyc}} \frac{a}{\sqrt{a^2+8bc}}
  = \sum_{\text{cyc}} \frac{a^2}{a\sqrt{a^2+8bc}}
  \;\ge\; \frac{S^2}{\displaystyle\sum_{\text{cyc}} a\sqrt{a^2+8bc}}
  \;\eqqcolon\; \frac{S^2}{T},
\]
by \cref{cor:engel} applied with numerators $a$ and denominators
$b = a\sqrt{a^2+8bc}$.

\emph{Step 2 (Cauchy--Schwarz).} Apply \cref{thm:cs} to the vectors
$\big(\sqrt a\big)$ and $\big(\sqrt{a\,(a^2+8bc)}\big)$:
\[
  T^2
  = \left(\sum_{\text{cyc}} \sqrt{a} \cdot
      \sqrt{a(a^2+8bc)}\right)^{\!2}
  \le S \cdot \left(a^3+b^3+c^3 + 24abc\right).
\]

\emph{Step 3 (expansion + AM--GM).} We claim
\[
  a^3+b^3+c^3+24abc \;\le\; S^3 .
\]
Expanding,
\[
  S^3 = a^3+b^3+c^3
  + 3\!\!\sum_{\text{sym}}\! a^2b \; + 6abc,
\]
where $\sum_{\text{sym}} a^2b = a^2b + a^2c + ab^2 + b^2c + ac^2 +
bc^2$. It remains to show $18abc \le 3\sum_{\text{sym}} a^2b$, i.e.\
$6abc \le \sum_{\text{sym}} a^2b$. This is AM--GM
(\cref{thm:amgm}) applied to the \emph{six} terms of the symmetric
sum: their product is
\[
  (a^2b)(a^2c)(ab^2)(b^2c)(ac^2)(bc^2) = a^6 b^6 c^6,
\]
so their geometric mean is $abc$ and their sum is at least $6abc$.

\emph{Conclusion.} Combining,
\[
  T \le \sqrt{S \cdot S^3} = S^2
  \qquad\Longrightarrow\qquad
  \sum_{\text{cyc}} \frac{a}{\sqrt{a^2+8bc}}
  \;\ge\; \frac{S^2}{T} \ge \frac{S^2}{S^2} = 1 .
\]
Equality in every step requires $a=b=c$: then each term equals
$\tfrac{a}{\sqrt{9a^2}} = \tfrac13$ and the sum is exactly $1$.
\end{solution}

\fi

% components/chapter-summary.tex — End-of-chapter distillation
% Reusable pattern: a chaptersummary box with \summaryitem bullets.
\begin{chaptersummary}
\begin{itemize}
  \summaryitem \textbf{AM--GM (\cref{thm:amgm2}, \cref{thm:amgm}).}
    $\dfrac{a_1 + \cdots + a_n}{n} \ge \sqrt[n]{a_1 \cdots a_n}$ for
    positive reals; equality iff all entries are equal. Proved by
    Cauchy's forward--backward induction: $P(2)$, $P(n) \Rightarrow
    P(2n)$, $P(n) \Rightarrow P(n{-}1)$.
  \summaryitem \textbf{Cauchy--Schwarz (\cref{thm:cs}).}
    $\big(\sum a_kb_k\big)^2 \le \big(\sum a_k^2\big)\big(\sum
    b_k^2\big)$; equality iff the vectors are proportional. Two
    proofs: the discriminant trick on $\sum (a_kt + b_k)^2 \ge 0$,
    and the Lagrange identity
    $\sum_{i<j}(a_ib_j - a_jb_i)^2$.
  \summaryitem \textbf{Engel form (\cref{cor:engel}).}
    $\sum \dfrac{a_k^2}{b_k} \ge \dfrac{(\sum a_k)^2}{\sum b_k}$ for
    $b_k > 0$: the homogenizing workhorse for fraction sums
    (\cref{str:engel}).
  \summaryitem \textbf{The equality case is the problem.} Every
    application above pins down \emph{when} equality holds; a
    solution without this analysis is incomplete.
\end{itemize}
\end{chaptersummary}

\begin{notation}
\begin{tabular}{@{}ll@{}}
  $\Positive$ & strictly positive reals \\
  $\sumk$ & $\sum_{k=1}^{n}$ \\
  $\sum_{\text{sym}} a^2b$ & $a^2b + a^2c + ab^2 + b^2c + ac^2 + bc^2$ \\
  $\sum_{\text{cyc}} f(a,b,c)$ & $f(a,b,c) + f(b,c,a) + f(c,a,b)$ \\
\end{tabular}
\end{notation}

